\section{Introduction}

%\jnote{I started writing some intro, it is quite wordy, maybe we should get rid of most of it, to get our point across much more strongly.}

Number on the forehead (NOF) communication complexity was introduced by Chandra, Furst and Lipton \cite{CFL83} as a variant of Yao's model of communication complexity. Here, we consider $k$ players, each of them having one of $k$ inputs, but instead of providing the players with their inputs, we imagine them having their inputs written on their foreheads --- in particular each player can observe all inputs \emph{except} for their own. Similar to the standard communication complexity, the players wish to exchange as little communication as possible, in order to jointly compute some function of their inputs.

As it turns out this innocuous alteration has deep mathematical consequences. For example, being able to prove lower bounds for specific NOF problems with $k = \Theta(\log^c n)$ players would imply explicit circuit complexity lower bounds \cite{HG90,BT94}.
%\jnote{I made this up, I vaguely remember Toni saying something about it at some point, and reading a statement to this effect in one of her papers... We need to check it.}
%\joshnote{I think this is Observation 17 here \url{https://www.cs.toronto.edu/~toni/Courses/CommComplexity2022/Lectures/notes-nof.pdf}}
%\jnote{Yes, thanks, that's what I had in mind.}

Moreover, even if we consider only three players, the theory of NOF communication complexity is simultaneously extremely deep, and frustratingly lacking. Its depth comes in part from a number of fascinating and perhaps surprising connections between natural questions in NOF communication complexity and well-studied problems in Ramsey Theory, extremal combinatorics and additive combinatorics. For example, already in the seminal paper \cite{CFL83}, a surprisingly efficient deterministic protocol  was provided for the NOF problem of checking whether $x+y+z = 0$ in $\bZ_N$, where $x,y,z \in \bZ_N$ are the inputs on the players' foreheads (we will call this problem $\eval_{\bZ_N}$ from now on). The protocol uses just $\Oh(\sqrt{\log N})$ communication, and makes use of Behrand's construction of a relatively dense set without three-terms arithmetic progression \cite{behrend}. On the other hand, in the same paper they showed an $\omega(1)$ lower bound for the $\eval_{\bZ_N}$ problem using the Hales–Jewett theorem --- a deep theorem in Ramsey theory. Since then, connections with the Ruzsa–Szemerédi problem, Corner-Free Sets and other objects from extremal combinatorics have been discovered, frequently leading to state-of-the-art constructions~\cite{LPS17, AS20,AMS12,LS21, CFTZ21}.

At the same time, perhaps because of these connections, many seemingly simple questions about NOF communication complexity remain unresolved, despite the fact that their number-in-hand variants are very easy to answer. For example, completely resolving the communication complexity of $\eval_{\bZ_N}$ requires a significantly better understanding of the correct quantitative dependency in the Roth's theorem (that large density subsets of $[N]$ contain a 3-terms arithmetic progression), and seems to be far out of reach of current techniques. 

Similarly, it is still open how to find an explicit $3$-party NOF problem which can be solved with $\Oh(1)$ randomized communication complexity (assuming shared randomness), but  which requires $\Omega(n)$ deterministic communication, even though this is a relatively simple property of $\textsc{Equality}$ for two-party communication, and even though, with a non-constructive argument, it is possible to show that such NOF problems exist \cite{DP08}. A natural candidate for such a problem is checking whether $x+y+z = 0$ in $\F_3^n$ (or, more generally, if $xyz = 1$ in a large power $G^n$ of some fixed group $G$, not necessarily abelian --- we will call those problems $\eval_{G^n}$), but proving a NOF communication complexity lower bound has been so far elusive.

Problems like $\eval_{G^n}$, in which for any pair of inputs to two of the players, there is exactly one input to the remaining player for which they should accept, will be called \emph{permutation problems}. These problems always have $\Oh(1)$ randomized communication complexity, so they are natural candidates to try to prove $\Omega(n)$ deterministic communication lower bounds. Indeed, a number of recent works have studied permutation problems toward this goal~\cite{BDPW10,CFTZ21}.

The construction \cite{CFL83} of an efficient protocol for $\eval_{\bZ_N}$ can be interpreted as a statement that existence of a large set $S \subset G$ without three terms arithmetic progressions in an abelian group $G$ can be used to provide an efficient protocol for $\eval_{G}$; by contrapositive, a lower bound for the communication complexity of $\eval_G$ would imply an upper bound for the size of any such set $S$. This latter statement for the group $\F_3^n$ used to be known as a \emph{cap-set conjecture} for few decades, and has famously been positively resolved in 2016 \cite{ED17,CLP17, T16} --- the largest set $S \subset \F_3^n$ without three terms arithmetic progressions has size $\Oh(c^n)$ for some $c < 3$. This leads to a natural question: can the techniques used to prove the cap-set theorem be generalized to prove a linear lower bound for the communication complexity of $\eval_{\F_3^n}$?

 Christandl et al. \cite{CFTZ21} recently identified a barrier against a direct applications of such techniques --- they used a method called \emph{combinatorial degenration} to prove that a \emph{subrank} of a specific tensor associated with the communication problem $\eval_{G^n}$ for an abelian group $G$ is near-maximal --- whereas an upper bound on this subrank would imply a desired lower-bound for communication complexity.

%\jnote{TODO: Introducing permutation problems, asymptotic communication complexity, matrix multiplication tensor and notions of subranks --- i.e. necessary terminology to state our results. Ideally not too formally.}


\subsection{Our contributions}
In this paper, we study three-player communication through the lens of \emph{promise number-in-hand problems}. For an alphabet $\Sigma$, a promise problem $(I,P)$ consists of two subsets $I \subset P \subset \Sigma \times \Sigma \times \Sigma$. In this problem, the three players are given as input $a,b,c \in \Sigma$, respectively, with the promise that $(a,b,c) \in P$, and their goal is to determine whether $(a,b,c) \in I$. For instance, usual number-in-hand communication corresponds to this model where $P = \Sigma \times \Sigma \times \Sigma$.

NOF problems can also be viewed as promise number-in-hand problems. Each player is given as an input a pair $(a_i, b_i) \in \Sigma^2$ for $i \in \{1, 2,3\}$, and their inputs are guaranteed to satisfy that $a_2 = b_1, a_3 = b_2$ and $a_1 = b_3$. (For instance, one thinks of $a_2$ and $b_1$ as the input on the head of player 3, so we are promised that players 1 and 2 see the same value there.) Write $\mmP \subset \Sigma^2 \times \Sigma^2 \times \Sigma^2$ to denote the inputs that satisfy that promise.

A key idea we pursue is to interpret $I$ and $P$, not just as subsets of inputs, but also as order-3 (3-dimensional) $\{0,1\}$-valued tensors. For example, letting $N = |\Sigma|$, we can represent $\mmP$ as a $\{0,1\}$-valued tensor $\mmP \in \F^{N^2 \times N^2 \times N^2}$, with the tensor entry being one whenever the corresponding tuple is a legal input.

As it turns out, $\mmP$ is \emph{exactly} the matrix multiplication tensor, the same tensor which is extensively studied in the theory of fast matrix multiplication algorithms! We leverage this extremely surprising, if only formal connection, to transfer ideas which have been developed in studying fast matrix multiplication to give new insights into NOF communication complexity, and conversely, to use results proved originally about NOF communication complexity and obtain new alternative proofs of consequential properties of the matrix multiplication tensor.

We use this connection to prove the following results.

    \paragraph{Communication Protocols from the Laser Method} We use the Laser method, the technique introduced by Strassen and used to design the best-known matrix multiplication algorithms~\cite{S86,CW87}, as a tool to design NOF communication protocols. For any three NOF problems $I, J, K$, let $I \otimes J \otimes K$ (``their product'') denote the NOF problem of solving the three simultaneously, i.e., given one input to each, determining whether all three are accepting inputs. Using the Laser method, we show that for any three permutation NOF problems $I, J, K$ with $n$-bit inputs, the product $I \otimes J \otimes K$ (which has $3n$-bit inputs) can be solved using only $n$ bits of communication deterministically.
    
    For example, for any permutation NOF problem $I$ with $n$-bit inputs, one can consider the problem of, given $N$ simultaneous inputs to $I$, determining if they are all accepting inputs. A consequence of our protocol is that this can be solved with $\frac13 n N + \Oh(n)$ bits of communication deterministically. This implies constructions of surprisingly large corner-free sets for powers of arbitrary abelian groups $G$.
    
    By contrast, we show using a counting argument that a random NOF permutation problem $I$ with $n$-bit inputs has requires deterministic communication complexity $\geq n/3 - \Oh(1)$. It was not previously clear to what extent this is tight; our protocol shows that this bound can at least be achieved for any permutation problem which is itself the product of three permutation problems.
    
    Our protocol should also be contrasted with the proof by probabilistic method in \cite{CFTZ21}, which showed that every NOF permutation problem with $n$-bit inputs has a protocol with communication complexity $n/2 + \Oh(1)$, improving in turn upon the trivial $n$ upper bound.
    %\jnote{We didn't define a permutation problem yet, nor product, nor asymptotic communication complexity...}
    %\jnote{What about corner-free sets for non-abelian groups? Is this a thing? Are they equivalent to protocols?}
    
     %In contrast, we show that random NOF permutation problems over the alphabet $[2^n] \times [2^n] \times [2^n]$ requires at least $n/3 - \Oh(1)$ communication.
    
    \paragraph{The Zeroing Out Subrank of Matrix Multiplication}
    A quantity which appears frequently in the theory of matrix multiplication algorithms is the \emph{zeroing out subrank} of matrix multiplication, denoted $Q_{zo}(\langle n,n,n \rangle)$. This is the size of the largest identity tensor (3-dimensional analogue of the identity matrix) which can be achieved as a zeroing out of the tensor $\langle n,n,n \rangle$ for multiplying two $n \times n$ matrices. Strassen~\cite{S86} famously proved that $Q_{zo}(\langle n,n,n \rangle) \geq n^{2 - o(1)}$, which is roughly as large as possible (since $\langle n,n,n \rangle$ is an $n^2 \times n^2 \times n^2$ tensor), and used it as part of the Laser method to design a faster algorithm for matrix multiplication\footnote{See, e.g.,~\cite[Section 8]{B13}. Bounds used in algorithms are sometimes described as bounds on a `monomial degeneration subrank' or `combinatorial degeneration subrank', but these are in turn ultimately used to bound the zeroing out subrank via~\cite{B80}. Note that bounds on more general notions like `border subrank' or `subrank' would not suffice in these algorithms, since the subrank decomposition is applied only to the `outer structure' of a tensor in the Laser method~\cite{S86}.}. Recent \emph{barrier} results have also used this bound to rule out certain approaches to designing faster matrix multiplication algorithms~\cite{AW21,A21,CVZ21}.

    In the bound $Q_{zo}(\langle n,n,n \rangle) \geq n^{2 - o(1)}$, it is natural to ask what is hidden by the $o(1)$. For instance, it was recently shown that the `border subrank' of $\langle n,n,n \rangle$ is $\lceil \frac34 n^2 \rceil$~\cite{KMZ20}; is it possible that we also have $Q_{zo}(\langle n,n,n \rangle) \geq \Omega(n^2)$? This could improve low-order terms in some applications.
    
    In fact, we prove that this is not possible, and that $Q_{zo}(\langle n,n,n \rangle) < n^2 / \omega(1)$. To prove this, we build off of the work of~\cite{LPS17}, who showed a NOF communication lower bound, that any NOF permutation problem requires $\omega(1)$ deterministic communication. Using our connection between NOF communication and matrix multiplication, we are able to modify their proof to get our upper bound on $Q_{zo}(\langle n,n,n \rangle)$.
    
    Interestingly, the prior work~\cite{LPS17} proved their communication lower bound using a connection with the Ruzsa–Szemerédi problem from extremal combinatorics. Not only does our new proof show a connection between $Q_{zo}(\langle n,n,n \rangle)$ and the Ruzsa–Szemerédi problem, but we are in fact able to show that the two are \emph{equivalent}! More precisely, it is impossible to determine the precise asymptotics of $Q_{zo}(\langle n,n,n \rangle)$ without also resolving the Ruzsa–Szemerédi problem. This is interesting in contrast with the `border subrank' which has been precisely determined~\cite{KMZ20}.
    
     Finally, using the connection with NOF communication, we also provide an alternative proof of Strassen's lower bound $Q_{zo}(\langle n,n,n \rangle) \geq n^{2 - o(1)}$, deducing it as a corollary of the existence of the efficient deterministic NOF protocol for $\eval_{\bZ_N}$.

     \paragraph{Slice-Rank Methods}

     Since their introduction~\cite{CFL83}, NOF problems $I$ have been associated with a corresponding 3-hypergraph $H$ whose chromatic number exactly characterizes the deterministic communication complexity of $I$. One recent powerful approach for bounding the chromatic number of such a hypergraph is the `slice rank' method, which was introduced to resolve the cap-set conjecture~\cite{CLP17,ED17}. Roughly speaking, if one can show that the slice rank of the adjacency tensor of $H$ is not too large, then this implies the chromatic number of $H$ is not too small, and hence a communication lower bound.
     
     Since then, a few successful techniques have been introduced to bound the slice ranks of many different tensors~\cite{T16,BCHGNSU17}. Recent work~\cite{CFTZ22} asked whether these techniques could give linear deterministic NOF communication lower bounds. However, they showed that, unfortunately, it is impossible to use such techniques to prove communication lower bounds for any $\textsc{Eval}$ group problem. They proved this by defining the corresponding adjacency tensor and calculating that it has maximal `asymptotic subrank' and hence large slice rank.

     We generalize their result and show that slice-rank methods cannot be used to prove a linear deterministic communication lower bound for \emph{any} NOF permutation problem (not just an $\textsc{Eval}$ group problem). To show this, we observe that for any NOF permutation problem, the corresponding adjacency tensor is actually (a permutation of) the matrix multiplication tensor $\mmP$. Since the matrix multiplication tensor is known to have maximal asymptotic subrank~\cite{S86}, it follows that the adjacency tensor never has small enough slice rank. In particular, this implies that in hindsight, the tensor whose asymptotic subrank was computed in the prior work~\cite{CFTZ22} is exactly the matrix multiplication tensor.

     By contrast, we use slice-rank methods to prove linear lower bounds for other promise problems. The key is that this method works whenever the \emph{promise tensor} has low slice-rank. For example, consider the promise problem $\eq_{\F_3^N}$, where parties receive inputs $x,y,z \in \F_3^N$ with a promise that $x+y+z = 0$, and they wish to decide whether $x=y$. We prove that this problem requires $\Omega(N)$ communication. This may seem like it should be straightforward to prove (``how could the third player's input $z=-(x+y)$ possibly help the first two players to test for equality?''), but we prove that the problem is actually \emph{equivalent} to the cap-set problem (which was a long-open conjecture).

     \paragraph{General Rank Methods and the Asymptotic Spectrum of Tensors}

     Slice-rank is one way to measure the `complexity' of a tensor, but there are many others, including other notions of tensor rank, and more generally a spectrum of different measures which is well-studied in algebraic complexity and tensor combinatorics called the `asymptotic spectrum of tensors'~\cite{S86,WZ22}. Let $r$ be any of these measures (the reader may want to focus on the case when $r$ is the tensor rank). We generalize the slice-rank method and prove that for any promise problem $(I,P)$, its deterministic communication complexity is at least $\log_2(r(I) / r(P))$. Hence, finding a measure $r$ which is larger for the problem $I$ than the promise $P$ leads to a communication lower bound. (This is why, for slice rank, the fact that the promise $\mmP$ has maximal slice rank means that a lower bound cannot be proved in this way.)

     This already has intriguing consequences when $r$ is tensor rank $R$.
     
     For number-in-hand communication, we know that the promise tensor $P_{NIH}$ is the all-1s tensor which has rank $R(P_{NIH}) = 1$. Hence, for any number-in-hand problem $I$, we get a communication lower bound of $\log_2(R(I))$. This is analogous to the standard upper bound in two-party deterministic communication compleixty by the log of the rank of the communication matrix. Moreover, any $2^n \times 2^n \times 2^n$ tensor $I$ (corresponding to an $n$-bit input problem) which is not `degenerate' in some way has rank at least $2^n$, which explains why nearly all number-in-hand communication problems require linear communication. 

     For NOF communication for $n$-bit inputs, as we've discussed, the promise tensor $P_{NOF}$ is the matrix multiplication tensor. Determining the rank of matrix multiplication is the \emph{central question} when designing matrix multiplication algorithms; the matrix multiplication exponent $\omega$ is defined exactly such that the rank of $P_{NOF}$ is $R(P_{NOF}) = (2^n)^{\omega + o(1)}$. Hence, for any problem $I$ with rank $R(I) \geq (2^n)^c$, we get a communication lower bound of $(c - \omega - o(1)) \cdot n$. The fact that $\omega$ appears negated in this lower bound is intriguing: this means that designing faster matrix multiplication algorithms, and hence lowering the known upper bound on $\omega$, leads to an improved communication lower bound for such problems $I$. 

     \paragraph{Matrix Multiplication as a Communication Problem}

     Finally, we consider how other promise problems can be used to shed further light on NOF communication lower bounds. Let $I$ be any $n$-bit NOF problem, let $N = 2^n$, and let $T_{(\mathbb{Z} / N\mathbb{Z})^2}$ denote the structure tensor of the group $(\mathbb{Z} / N\mathbb{Z})^2$, i.e., the tensor
$$T_{(\mathbb{Z} / N\mathbb{Z})^2} = \sum_{a,b,c,d \in [N]} x_{(a,b)} y_{(c,d)} z_{(a+c\pmod{N}, b+d\pmod{N})}.$$
    We observe that $P_{NOF} \subset T_{(\mathbb{Z} / N\mathbb{Z})^2}$ (after appropriately permuting variables), and so it is well-defined to consider three different promise problems using these three tensors: $(I,P_{NOF})$, $(P_{NOF},T_{(\mathbb{Z} / N\mathbb{Z})^2})$, and $(I,T_{(\mathbb{Z} / N\mathbb{Z})^2})$. Moreover, a simple triangle inequality-style argument shows that their communication complexities are related via:
    $$\cc(I,P_{NOF}) \geq \cc(I,T_{(\mathbb{Z} / N\mathbb{Z})^2}) - \cc(P_{NOF},T_{(\mathbb{Z} / N\mathbb{Z})^2}).$$
    On the left-hand side is exactly the deterministic NOF complexity of the problem $I$ that we would like to prove a lower bound for. On the right-hand side, $(I,T_{(\mathbb{Z} / N\mathbb{Z})^2})$ is a similar problem to the original NOF problem of $I$, but with a weaker promise, and so it should be easier to prove communication lower bounds for. What about $\cc(P_{NOF},T_{(\mathbb{Z} / N\mathbb{Z})^2})$? (Recall that in this problem, we're given as input a term from the tensor $T_{(\mathbb{Z} / N\mathbb{Z})^2}$, and our goal is to determine whether or not it is a valid NOF input from $P_{NOF}$.)

    For this problem, we give a protocol with communication $\leq n$ (the trivial bound would be $2n$). Moreover, we prove a communication lower bound of $\geq (\omega - 2) n$ (this follows from the rank-based approach discussed above). So, further improvements to the communication protocol actually \emph{require} fast matrix multiplication (since they would imply an upper bound on $\omega$ via this inequality). If it were the case that $\omega = 2$ (which is popularly conjectured), then it is plausible that this problem actually has sublinear communication complexity. In that case, in order to prove a linear lower bound for the NOF problem $(I,P_{NOF})$, it would suffice to prove a linear lower bound for the seemingly-harder problem $(I,T_{(\mathbb{Z} / N\mathbb{Z})^2})$. In summary, if we could design a faster matrix multiplication algorithm achieving $\omega=2$, then this may be a promising approach to proving NOF communication lower bounds.